\documentclass[12pt]{amsart}
%   \overfullrule=5pt
\usepackage[active]{srcltx}
\usepackage{calc,amssymb,amsthm,amsmath,amscd, eucal,ulem}
\usepackage{alltt}
%\usepackage[left=1in,top=1in,right=1in,bottom=1in]{geometry}
\synctex=1
%\usepackage{mathtools}
\RequirePackage[dvipsnames,usenames]{color}
\let\mffont=\sf
\normalem
\input{mabliautoref.sty}
\input{xy}
\xyoption{all}
\input{kmacros3.sty}
\usepackage{tikz}
\usepackage{amsfonts, mathrsfs}
\usepackage{calligra}
\usepackage{stmaryrd} %power series bracket
%\usepackage{dcpic, pictexwd}
%\usepackage{pdfsync}
\usepackage[left=1.1in,top=1in,right=1.1in,bottom=1in]{geometry}
\usepackage{enumitem}

\numberwithin{equation}{theorem}

%\renewcommand{\thefootnote}{\fnsymbol{footnote}}
\newcommand{\D}{\displaystyle}
\newcommand{\til}{\widetilde}
\newcommand{\ol}{\overline}
\newcommand{\F}{\mathbb{F}}
\DeclareMathOperator{\E}{E}
\renewcommand{\:}{\colon}
\newcommand{\eg}{{\itshape e.g.} }
\renewcommand{\m}{\mathfrak{m}}
\renewcommand{\n}{\mathfrak{n}}
\newcommand{\Tt}{{\mathfrak{T}}}
\newcommand{\calC}{\mathcal{C}}
\newcommand{\RamiT}{\mathcal{R_{T}}}
\DeclareMathOperator{\edim}{edim}
\DeclareMathOperator{\length}{length}
\DeclareMathOperator{\monomials}{monomials}
\DeclareMathOperator{\Fitt}{Fitt}
\DeclareMathOperator{\depth}{depth}
\newcommand{\Frob}[2]{{#1}^{1/p^{#2}}}
\newcommand{\Frobp}[2]{{(#1)}^{1/p^{#2}}}
\newcommand{\FrobP}[2]{{\left(#1\right)}^{1/p^{#2}}}
\newcommand{\extends}{extends over $\Tt${}}
\newcommand{\extension}{extension over $\Tt${}}
\newcommand{\fg}{\pi_1^{\textnormal{\'{e}t}}} %Fundamental group
\DeclareMathOperator{\ns}{ns}
%\newcommand{\Spec}{\text{Spec}} %Spectrum
\newcommand{\pSpec}{\text{Spec}^{\circ}} %Puntured Spectrum
\newcommand{\etInCdOne}{\etale{} in codimension $1$}
\DeclareMathOperator{\DIV}{Div}
%\renewcommand{\char}{\charact}


\usepackage{setspace}
\usepackage{hyperref}
%\usepackage{enumerate}
\usepackage{graphicx}
\usepackage[all,cmtip]{xy}

\usepackage{verbatim}

\theoremstyle{theorem}
\newtheorem*{mainthm}{Main Theorem}
\newtheorem*{mainthma}{Main Theorem A}
\newtheorem*{mainthmb}{Main Theorem B}
\newtheorem*{mainthmc}{Main Theorem C}
\newtheorem*{atheorem}{Theorem}


\newcommand{\sol}[1]{ \vskip 6pt{\color{blue} {\bf Solution:} #1} \vskip 6pt}
%The todo box!
\def\todo#1{\textcolor{Mahogany}%
{\footnotesize\newline{\color{Mahogany}\fbox{\parbox{\textwidth-15pt}{\textbf{todo:
} #1}}}\newline}}

%What is going on?  Why isn't \O a script O?!?
\renewcommand{\O}{\mathscr O}

\begin{document}
\title{Exercises for Characteristic $p$ commutative algebra\\February 27rd, 2017}
\author{Due, Friday March 10th, 2017}
%\address{Department of Mathematics\\ University of Utah\\ Salt Lake City\\ UT 84112}
%\email{schwede@math.utah.edu}


%  \subjclass[2010]{14F18, 13A35, 14B05}

\maketitle

%\section{Reduction to characteristic $p > 0$}
%\chapter{Frobenius and Kunz's theorem}
%\bibliographystyle{skalpha}
%\bibliography{MainBib}

\begin{enumerate}[label=(\arabic*)]
\item For $R$ a Noetherian ring, $I \subseteq R$ an ideal and $U = \Spec R \setminus V(I)$ we defined for any $R$-module $M$
\[
\Gamma(U, M) = \ker \Big( \bigoplus_j M_{f_j} \to \bigoplus_{a < b} M_{f_{a} f_b} \big)
\]
where $I = \langle f_1, \ldots, f_m \rangle$.  Prove that $\Gamma(U, M)$ is independent of the choice of generators of $I$.
\item
Suppose that $R$ is a Noetherian ring.  Recall that an object $\omega^{\mydot} \in D^b_{f.g.}(R)$ is called a \emph{dualizing complex for $R$} if the following two conditions are satisfied.
\begin{itemize}
\item[(a)] $\omega^{\mydot}$ has finite injective dimension (is quasi-isomorphic to a bounded complex of injectives) and,
\item[(b)] The functor $\myD(\blank) = \myR \Hom_R(\blank, \omega^{\mydot})$ has the property that the canonical map $C^{\mydot} \to \myD(\myD(C^{\mydot}))$ is an isomorphism for all $C^{\mydot} \in D^b_{f.g.}(R)$.
\item[(b')] Or equivalently to (b), $R \cong \myR \Hom_R(\omega^{\mydot}, \omega^{\mydot})$.
\end{itemize}
Prove that (b') implies (b) above.
\vskip 3pt
\emph{Hint:} If you are stuck you may use the following formula (or prove it if you are feeling ambitious).  For $A^{\mydot}, B^{\mydot}, C^{\mydot} \in D^b_{f.g.}(R)$ with $C^{\mydot}$ quasi-isomorphic to a finite complex of injectives, we have
\[
A^{\mydot} \otimes^{\myL}_R \myR \Hom_R(B^{\mydot}, C^{\mydot}) \qis \myR \Hom_R(\myR \Hom_R(A^{\mydot}, B^{\mydot}), C^{\mydot})
\]
\item Suppose that $(R, \fram)$ is an $F$-finite Noetherian local ring of characteristic $p > 0$ with a dualizing complex.  Then $\myR \Hom_R(F^e_* R, \omega_R^{\mydot}) =: \omega_{F^e_* R}^{\mydot}$ is a dualizing complex for $F^e_* R$.  Furthermore, $F^e_* \omega_R^{\mydot} \qis \omega_{F^e_* R}^{\mydot}$.
    \vskip 3pt
    \emph{Hint:} The fact that it is a dualizing complex was worked out in class.  Dualizing complexes are unique up to shifting and twisting by rank-one projectives.  Use the fact that it is local to handle the rank-1 projective case.  Then use the that the localization of a dualizing complex is a dualizing complex to handle the shift (localize at a minimal prime).
\item Show that the formation of $R^{\myN}$ (the normalization of $R$) commutes with localization, in particular if $W \subseteq R$ is a multiplicative set then $(W^{-1} R)^{\myN} = W^{-1} (R^{\myN})$.
\item Suppose that $R = k[x]$ ($\mathrm{char} k \neq 2$) and consider the three points $x = -1, x = 0, x = 1$.  Glue these three points together and find a presentation of the resulting non-normal ring.  Give a geometric justification about why the $\Spec$ of it cannot be embedded as a closed subscheme of $\bA^2 = \Spec k[u,v]$.
\item Fix $R$ an $F$-finite ring of characteristic $p > 0$, $\fra \subseteq R$ an ideal and $t \geq 0$ a real number.  We say that a pair $(R, \fra^t)$ is \emph{(sharply) $F$-split} if there exists an $e > 0$ and an element
\[
\phi \in (F^e_* \fra^{\lceil t(p^e - 1) \rceil}) \cdot \Hom_R(F^e_* R, R)
\]
such that $\phi(F^e_* R) = R$.  Prove the following facts about sharply $F$-split pairs.
\begin{enumerate}
\item Show that if $(R, \fra^t)$ is sharply $F$-split then for all sufficiently divisible $e > 0$, there exists a $\phi$ as above.
\item Find an example that is sharply $F$-split but such that a $\phi$ as above does not exist for every $e > 0$.
\item Suppose that $(R, \fra^t)$ is sharply $F$-split.  Show that there exists $t' \geq t$ with $t' = {b \over p^e - 1}$ such that $(R, \fra^{t'})$ is also sharply $F$-split.
\item Show that $(R, \fra^t)$ is sharply $F$-split if and only if for every maximal $\fram \in \Spec R$, $(R_{\fram}, \fra_{\fram}^t)$ is sharply $F$-split.
\item If $(R, \fram)$ is local, then show that $(R, \fra^t)$ is sharply $F$-split if and only if for some $e > 0$ there exists an element $z \in \fra^{\lceil t(p^e - 1)\rceil}$ such that the map
\[
\xymatrix@R=0pt{
R \ar[r] & F^e_* R \\
1 \ar@{|->}[r] & F^e_* z \\
}
\]
splits as a map of $R$-modules.
\end{enumerate}
\item A map of rings $R \to S$ is called \emph{pure} if for every $R$-module $M$, the map $M \to M \otimes_R S$ is injective.  Suppose that $S$ is a finite $R$-module, show that $R \to S$ is pure if and only if $R \to S$ splits as a map of $R$-modules.
    \vskip 3pt
    \emph{Hint:} First reduce to the case that $R$ is complete and local and let $E$ be the injective hull of the residue field.  Consider the injective map $E \to E \otimes_R S$ and apply Matlis duality and $\Hom-\tensor$ adjointness.
    \vskip 3pt
    \noindent
    {\bf Remark.} A ring is called \emph{$F$-pure} if the Frobenius map $R \to F_* R$ is pure.  Because of this, and because most rings we study are $F$-finite, the terms $F$-pure or $F$-split are usually synonymous.
\item Given an ideal $\fra \subseteq R$ in an $F$-finite ring, we define the \emph{$F$-split (or $F$-pure) threshold $\fpt(R, \fra)$ of $(R, \fra)$} to be
\[
\fpt(R, \fra) := \{ t \geq 0\;|\; (R, \fra^t) \text{ is sharply $F$-split} \}.
\]
Compute the $F$-pure threshold of the following pairs, in all of them $k = \overline{\bF_p}$.
\begin{enumerate}
\item $k[x,y]$, $\fra = \langle x^u y ^v \rangle$.
\item $k[x,y,z]$, $\fra = \langle x^2 y - z^2 \rangle$ (the answer depends on the characteristic).
\item $k[x,y]$, $\fra = \langle y^2 - x^3 \rangle$ (the answer also depends on the characteristic).
\end{enumerate}
\end{enumerate}



\end{document}


